%========================================================================%
% Results section content. Slots into main.tex below the \section{Results}
% header and the Pareto figure block (Fig.~\ref{fig:pareto}).
%========================================================================%

%-----------------------------------------------------------------------%
\subsection{Multi-window question answering}\label{sec:squad}
%-----------------------------------------------------------------------%

The controlled multi-window QA task partitions a SQuAD validation example into five sequential windows: four fact windows containing Wikipedia paragraphs followed by a query window with the question. The answer-bearing window is rotated across the four non-query positions, so accuracy at each position isolates the cache-management problem rather than the model's recall bias toward window-final content. All methods receive the same pre-segmented windows and the same total cache budget $k{=}300$, so accuracy depends solely on which cached entries are retained and how they are augmented. We report position-robust accuracy, defined as the mean over the four answer-bearing positions of the per-position accuracy on $n{=}100$ samples per position.

\Cref{tab:squad_main} reports the five-seed mean$\pm$std (seeds 7, 42, 123, 999, 2024) of position-robust accuracy for Stage~1 (parameter-free multiplicative selection at $k{=}300$) and the full pipeline (16 learned summary tokens prepended to 284 multiplicatively selected real tokens, total budget $k{=}300$). For context, we also report KIVI~\cite{liu2024kivi} at K4V4 applied to the multiplicatively selected cache at $k{=}300$ from a single canonical seed, which we take as the strongest published cache-compression baseline against which our pipeline is positioned.

\begin{table}[t]
\caption{Position-robust multi-window SQuAD accuracy (\%). Stage~1 and Stage~1+2 columns are five-seed mean$\pm$std (seeds 7/42/123/999/2024, $n{=}100$ samples per seed, four answer-bearing positions per sample). KIVI K4V4 is the same multiplicatively selected cache at $k{=}300$ further compressed with 4-bit asymmetric per-channel keys and per-token values, single canonical seed. The Stage~1+2 pipeline holds the total budget fixed at $k{=}300$. For Mistral 7B and Mistral-Small 24B we report the longer-context retrain (n=10 windows, mixed SQuAD/HotpotQA); the four other backbones use the original n=5 SQuAD-only training.}
\label{tab:squad_main}
\centering
\small
\setlength{\tabcolsep}{4pt}
\begin{tabular}{l ccc}
\toprule
\textbf{Model} & \textbf{Stage 1 (\%)} & \textbf{Stage 1+2 (\%)} & \textbf{KIVI K4V4 (\%)} \\
\midrule
Qwen 2.5 7B   & $59.8 \pm 2.7$  & $66.8 \pm 4.0$  & 45.5 \\
Qwen 2.5 14B  & $63.8 \pm 5.3$  & $73.8 \pm 2.6$  & 68.0 \\
Qwen 2.5 32B  & $64.7 \pm 4.9$  & $72.5 \pm 2.9$  & 69.0 \\
Llama 3.1 8B  & $56.5 \pm 4.0$  & $65.9 \pm 4.5$  & 55.0 \\
Mistral 7B    & $44.1 \pm 1.8$  & $57.0 \pm 3.5$  & 36.0 \\
Mistral-Small 24B & $59.9 \pm 4.0$  & $63.6 \pm 3.3$  & 59.5 \\
\bottomrule
\end{tabular}
\end{table}

Across the six backbones, adding the learned summary module to Stage~1 selection produces a mean improvement of $+8.5$ percentage points (range $+3.6$ to $+12.9$), with non-overlapping seed envelopes on Qwen~14B, Qwen~32B, Llama~8B, and Mistral~7B. The gap between Stage~1+2 and KIVI K4V4 ranges from $+4.1$ pp on Mistral-Small~24B to $+21.0$ pp on Mistral~7B; on Qwen~7B and Llama~8B the gap exceeds $10$ pp despite KIVI consuming $\sim$2.5$\times$ more bits per token than the FP16 reference. Pushing KIVI to two-bit precision (K2V2) degrades it further: 12.0\% on Qwen~7B, 47.0\% on Llama~8B, 60.0\% on Qwen~14B, 32.5\% on Mistral~7B and 58.5\% on Mistral-Small~24B, all of which sit below Stage~1+2 by 14 to 67 percentage points. Position-robust evaluation is more stringent than a single fixed-position number: on positions far from the query (positions~2 and~3 in our setup), pure selection routinely loses information that the learned summary recovers. The mean Stage~1+2 accuracy increases monotonically through the Qwen family ($66.8 \to 73.8 \to 72.5$, with 14B and 32B within one standard deviation) and rises from Mistral~7B to Mistral-Small~24B by $+6.6$ pp.

The Pareto frontier in \Cref{fig:pareto} places these accuracy numbers against KV-cache memory per request. The Stage~1+2 pipeline at $k{=}300$ FP16 consumes 16--75~MB of KV memory across the six models versus 420~MB--2.0~GB for the full $n{=}20$-window cache; with the K8V4 quantization stack (\cref{sec:compression}), the same pipeline operates at $25.6\times$ FP16 compression of the full cache and reaches $68\times$ when accuracy and memory are jointly considered.


%-----------------------------------------------------------------------%
\subsection{Transfer to standard long-context benchmarks}\label{sec:longbench}
%-----------------------------------------------------------------------%

To test whether the same drop-in transfers beyond our training distribution, we evaluate on two LongBench~\cite{bai2024longbench} extractive QA tasks: HotpotQA (multi-hop reasoning across multiple Wikipedia paragraphs) and MultiFieldQA (single-hop QA over diverse-domain documents). LongBench documents are chunked into 384-token windows and processed sequentially through the same pipeline used for SQuAD; the Stage~1+2 checkpoints are the same SQuAD-trained models from \cref{sec:squad} with no LongBench-specific tuning. The cache budget is $k{=}300$ in all cases.

\begin{table}[t]
\caption{LongBench F1 scores at $k{=}300$ (100 samples per task). Hot.\ = HotpotQA, MF = MultiFieldQA. \textcolor{red}{\textbf{[PLACEHOLDER]}} Faithful single-prompt SnapKV/H$_2$O/PyramidKV baselines are being re-run with the upstream KVCache-Factory monkey-patches; results pending. The previous baseline columns in this table used unfaithful re-implementations that this revision withdraws (see AUDIT\_BULLSHIT.md). Streaming-protocol baselines (matching AstroNet's operating regime) are under construction. AstroNet S1+S2 was trained on SQuAD/HotpotQA mixture with no LongBench-train exposure; absolute scores are low across all methods at the $k{=}300$ budget on multi-thousand-token documents.}
\label{tab:longbench}
\centering
\small
\setlength{\tabcolsep}{2.5pt}
\begin{tabular}{l cc cc cc cc cc cc}
\toprule
& \multicolumn{2}{c}{\textbf{Qwen 7B}} & \multicolumn{2}{c}{\textbf{Qwen 14B}} & \multicolumn{2}{c}{\textbf{Qwen 32B}} & \multicolumn{2}{c}{\textbf{Llama 8B}} & \multicolumn{2}{c}{\textbf{Mist.\ 7B}} & \multicolumn{2}{c}{\textbf{Mist.\ 24B}} \\
\cmidrule(lr){2-3} \cmidrule(lr){4-5} \cmidrule(lr){6-7} \cmidrule(lr){8-9} \cmidrule(lr){10-11} \cmidrule(lr){12-13}
& Hot. & MF & Hot. & MF & Hot. & MF & Hot. & MF & Hot. & MF & Hot. & MF \\
\midrule
\multicolumn{13}{l}{\textcolor{red}{\textit{Single-prompt baselines (oracle ceiling) — faithful upstream impl, PENDING:}}} \\
SnapKV (single-prompt)    & \multicolumn{12}{c}{\textcolor{red}{[PENDING — queue running, ETA $\sim$24 h]}} \\
H${}_2$O (single-prompt)  & \multicolumn{12}{c}{\textcolor{red}{[PENDING — queue running]}} \\
PyramidKV (single-prompt) & \multicolumn{12}{c}{\textcolor{red}{[PENDING — queue running]}} \\
\midrule
\multicolumn{13}{l}{\textcolor{red}{\textit{Streaming-protocol comparison (matched operating regime) — PENDING:}}} \\
SnapKV (streaming)        & \multicolumn{12}{c}{\textcolor{red}{[PENDING — protocol design in progress]}} \\
H${}_2$O (streaming)      & \multicolumn{12}{c}{\textcolor{red}{[PENDING]}} \\
\midrule
AstroNet S1     & 24.3 & 13.6 & 18.1 & 19.4 & 15.9 & 20.9 & 16.4 & 12.4 &  0.7 &  5.9 &  5.1 & 11.5 \\
AstroNet S1+S2  & 23.4 & 15.7 & 29.3 & 14.2 & 24.3 & 16.6 & 23.1 & 15.1 &  6.5 &  8.5 & 10.1 & 11.9 \\
\bottomrule
\end{tabular}
\end{table}

\textcolor{red}{\textbf{[PROSE PLACEHOLDER]}} Comparison narrative against published baselines will be rewritten once the faithful single-prompt and streaming-protocol baseline numbers land. The internal $\Delta$(S1$\to$S1+S2) signal is the load-bearing claim and is consistent in this table: adding the 16 learned summary tokens lifts HotpotQA F1 on Qwen 14B from 18.1 to 29.3 ($+11.2$) and Qwen 32B from 15.9 to 24.3 ($+8.4$); MultiFieldQA gains are smaller (1--3 F1). Absolute scores remain low across all methods at $k{=}300$ on multi-thousand-token documents. The Stage~1+2 checkpoints are trained on a SQuAD/HotpotQA mixture and never see LongBench MultiFieldQA documents, so MultiFieldQA probes transfer.


%-----------------------------------------------------------------------%
\subsection{Depth-robustness in needle-in-a-haystack}\label{sec:needle}
%-----------------------------------------------------------------------%

We test retrieval robustness by hiding a single answer-bearing fact among $n{-}1$ distractor windows at five depth positions (start, 25\%, 50\%, 75\%, end) with 20 trials per depth. At $n{=}5$ windows every method tested reaches near-perfect retrieval; we report $n{=}20$ in \cref{tab:needle}, the most stringent setting in our protocol, where the difference between cache-management strategies is most pronounced.

\textcolor{red}{\textbf{[PROSE WITHDRAWN]}} The previous version of this section claimed a ``faithful reimplementation'' of SnapKV. On audit, those baseline columns were hand-rolled approximations applied within AstroNet's protocol rather than faithful re-implementations of the published methods; see AUDIT\_BULLSHIT.md. The faithful single-prompt baselines are being re-run with the upstream KVCache-Factory monkey-patches; results pending.

\begin{table}[t]
\caption{Needle-in-a-haystack accuracy (\%) at $n{=}20$ windows, averaged over 5 depths $\times$ 20 trials, $k{=}300$. \textcolor{red}{\textbf{[PLACEHOLDER]}} Baseline columns withdrawn pending faithful re-runs. Internal $\Delta$(S1$\to$S1+S2) shows the load-bearing $+$19 to $+$50 pp lift from the 16 learned summary tokens (Qwen 32B and 14B saturate at 100\%). The Mistral-Small~24B Stage~1+2 entry uses the longer-context retrained checkpoint.}
\label{tab:needle}
\centering
\small
\setlength{\tabcolsep}{6pt}
\begin{tabular}{l c c c c c}
\toprule
\textbf{Model} & \textbf{StreamingLLM} & \textbf{H${}_2$O} & \textbf{SnapKV} & \textbf{AstroNet S1} & \textbf{AstroNet S1+S2} \\
\midrule
Qwen 2.5 7B       & \textcolor{red}{[PEND]} & \textcolor{red}{[PEND]} & \textcolor{red}{[PEND]} & 72 & \textbf{91} \\
Qwen 2.5 14B      & \textcolor{red}{[PEND]} & \textcolor{red}{[PEND]} & \textcolor{red}{[PEND]} & 79 & \textbf{100} \\
Qwen 2.5 32B      & \textcolor{red}{[PEND]} & \textcolor{red}{[PEND]} & \textcolor{red}{[PEND]} & 50 & \textbf{100} \\
Llama 3.1 8B      & \textcolor{red}{[PEND]} & \textcolor{red}{[PEND]} & \textcolor{red}{[PEND]} & 58 & \textbf{95} \\
Mistral 7B        & \textcolor{red}{[PEND]} & \textcolor{red}{[PEND]} & \textcolor{red}{[PEND]} & 5 & \textbf{32} \\
Mistral-Small 24B & \textcolor{red}{[PEND]} & \textcolor{red}{[PEND]} & \textcolor{red}{[PEND]} & 72 & \textbf{81} \\
\bottomrule
\end{tabular}
\end{table}

\textcolor{red}{\textbf{[PROSE PLACEHOLDER]}} Comparison narrative against StreamingLLM/H$_2$O/SnapKV will be rewritten once faithful baselines land. Internal observation: Stage~1+2 lifts AstroNet S1 by $+19$, $+21$, $+50$, and $+37$ percentage points on the four Qwen and Llama backbones at $n{=}20$, attributable to the 16 learned summary tokens. Mistral~7B shows a smaller $+27$ pp lift from a low base, consistent with weaker multi-hop capability of the base model at this cache budget. Faithful baseline numbers must be substituted before the headline ``beats SnapKV by $+X$~pp'' phrasing can be reintroduced.


%-----------------------------------------------------------------------%
\subsection{Long-context generalisation (RULER 8\,k--32\,k)}\label{sec:ruler}
%-----------------------------------------------------------------------%

To probe behaviour at native long-context sequence lengths, we evaluate on a needle-style retrieval task at $n=22$, $44$, and $85$ windows, corresponding to approximately $8\,000$, $16\,000$, and $32\,000$ tokens of input context, following the RULER protocol~\cite{hsieh2024ruler}. The Stage~2 module was trained at $n=10$ windows ($\sim 5\,000$ tokens), so $8\,$k tokens is in-distribution, $16\,$k is mildly out-of-distribution, and $32\,$k is substantially out-of-distribution for the learned component. The parameter-free Stage~1 selection, by contrast, has no training distribution and applies uniformly at all context lengths.

\begin{table}[t]
\caption{RULER long-context retrieval accuracy (\%) at $k{=}300$, averaged over 5 depth positions $\times$ 10 trials per depth. Context length grows from $\sim8$\,k tokens at $n=22$ to $\sim32$\,k tokens at $n=85$. Stage~2 training distribution is $n=10$ windows; $n=22$ is in-distribution, $n=85$ is OOD. \textcolor{red}{\textbf{[PLACEHOLDER]}} Baseline rows withdrawn; faithful single-prompt SnapKV/H$_2$O at the same budgets pending.}
\label{tab:ruler}
\centering
\small
\setlength{\tabcolsep}{6pt}
\begin{tabular}{l c c c c c c}
\toprule
& \multicolumn{3}{c}{\textbf{Qwen 2.5 14B}} & \multicolumn{3}{c}{\textbf{Llama 3.1 8B}} \\
\cmidrule(lr){2-4} \cmidrule(lr){5-7}
\textbf{Method} & 8\,k & 16\,k & 32\,k & 8\,k & 16\,k & 32\,k \\
\midrule
StreamingLLM    & \textcolor{red}{[P]} & \textcolor{red}{[P]} & \textcolor{red}{[P]} & \textcolor{red}{[P]} & \textcolor{red}{[P]} & \textcolor{red}{[P]} \\
H2O             & \textcolor{red}{[P]} & \textcolor{red}{[P]} & \textcolor{red}{[P]} & \textcolor{red}{[P]} & \textcolor{red}{[P]} & \textcolor{red}{[P]} \\
SnapKV          & \textcolor{red}{[P]} & \textcolor{red}{[P]} & \textcolor{red}{[P]} & \textcolor{red}{[P]} & \textcolor{red}{[P]} & \textcolor{red}{[P]} \\
AstroNet S1     & 72 & 58 & 56 & 48 & 28 &  8 \\
AstroNet S1+S2  & \textbf{100} & \textbf{66} & 52 & \textbf{88} & \textbf{50} & \textbf{14} \\
\bottomrule
\end{tabular}
\end{table}

\textcolor{red}{\textbf{[PROSE PLACEHOLDER]}} The internal $\Delta$(S1$\to$S1+S2) shows the S2 contribution in isolation: $+28$, $+8$, and $-4$ pp at $8/16/32$\,k on Qwen 14B, and $+40$, $+22$, $+6$ pp on Llama 8B. The negative delta at $32$\,k on Qwen 14B sits at $\sim$6$\times$ S2's training context length and is the OOD edge of the current S2 training distribution. Comparison narrative against faithful single-prompt SnapKV/H$_2$O at the same budgets will be reintroduced once those baselines land.


%-----------------------------------------------------------------------%
\subsection{KV-cache compression}\label{sec:compression}
%-----------------------------------------------------------------------%

The pipeline composes with downstream cache quantisation. For each of the six backbones, the Stage~1+2 FP16 cache of $k{=}300$ entries can be further compressed to $300$ tokens at 8-bit keys / 4-bit values (K8V4) using a classical per-(layer, head) Lloyd--Max scheme (RoPE-compatible) that we describe in \cref{sec:quant}. \Cref{tab:compression} reports both the memory geometry and the accuracy cost.

\begin{table}[t]
\caption{KV-cache compression and accuracy cost. Memory columns report bytes per request for the $k{=}300$ retained cache (FP16 and K8V4) versus the full $n{=}20$-window cache. Accuracy columns report \emph{single-seed, fixed answer position} SQuAD accuracy (\%) of multiplicative Stage~1 selection at FP16 vs.\ the same cache compressed to K8V4 with per-(layer, head) Lloyd--Max quantisation (RoPE-compatible; \cref{sec:quant}); the position-robust five-seed values for Stage~1 and Stage~1+2 are reported in \cref{tab:squad_main}. Vanilla TurboQuant (random rotation $+$ Lloyd--Max) collapses to $0\%$ on Qwen~7B due to destruction of rotary positional encoding and was not pursued on the other backbones (\cref{sec:quant}).}
\label{tab:compression}
\centering
\small
\setlength{\tabcolsep}{4pt}
\begin{tabular}{l rrr rr rr}
\toprule
& \multicolumn{3}{c}{\textbf{Memory (MB / request)}} & \multicolumn{2}{c}{\textbf{Compression}} & \multicolumn{2}{c}{\textbf{Accuracy (\%)}} \\
\cmidrule(lr){2-4}\cmidrule(lr){5-6}\cmidrule(lr){7-8}
\textbf{Model} & Full $n{=}20$ & $k{=}300$ FP16 & $k{=}300$ K8V4 & FP16 & K8V4 & FP16 & K8V4 \\
\midrule
Qwen 2.5 7B       & 420  & 16.4 &  6.2 & 25.6$\times$ & 68.3$\times$ & 74.5 & 77.5 \\
Qwen 2.5 14B      & 1440 & 56.3 & 21.1 & 25.6$\times$ & 68.3$\times$ & 75.0 & 69.5 \\
Qwen 2.5 32B      & 1920 & 75.0 & 28.1 & 25.6$\times$ & 68.3$\times$ & 78.5 & 77.0 \\
Llama 3.1 8B      &  960 & 37.5 & 14.1 & 25.6$\times$ & 68.3$\times$ & 78.5 & 78.5 \\
Mistral 7B        &  960 & 37.5 & 14.1 & 25.6$\times$ & 68.3$\times$ & 72.5 & 72.0 \\
Mistral-Small 24B & 1200 & 46.9 & 17.6 & 25.6$\times$ & 68.3$\times$ & 79.5 & 72.5 \\
\bottomrule
\end{tabular}
\end{table}

The selection stage alone reduces the cache by $25.6\times$ relative to the full $n{=}20$-window cache; appending the quantisation stack reaches $68.3\times$. On four of the six backbones the K8V4 accuracy cost is within $\pm 1.5$\,pp of the FP16 reference, and on Qwen~7B K8V4 in fact improves accuracy by $+3.0$\,pp (the asymmetric 8/4-bit precision suppresses small fluctuations that hurt the FP16 baseline on this model). The two-sided range across all six models is $-7.0$\,pp on Mistral-Small~24B to $+3.0$\,pp on Qwen~7B, comfortably bounded for a $2.6\times$ further compression of the retained cache. The asymmetry between key precision (8 bits) and value precision (4 bits) reflects the different roles the two play in the attention computation: keys appear inside an inner product with the query and are sensitive to small perturbations, while values are weight-averaged and tolerate coarser representations. The same K8V4 setting therefore applies identically across all six backbones in our experiments.

TurboQuant~\cite{zandieh2026turboquant} applies a random orthogonal rotation before per-coordinate Lloyd--Max quantisation; the rotation is its defining innovation, making coordinates approximately i.i.d.\ Gaussian so that scalar quantisation becomes near-optimal. On a backbone using rotary positional embeddings~\cite{su2024rope}, however, this rotation destroys the paired-coordinate structure that RoPE relies on and reduces accuracy to $0\%$ (verified on Qwen~7B at 4-bit precision). We therefore replace the rotation with per-(layer, head) Gaussianisation: each attention head is normalised independently to $\mathcal{N}(0,1)$ using its own running statistics, the standard Lloyd--Max scalar codebook~\cite{lloyd1982least,max1960quantizing} is applied per coordinate, and denormalisation is the inverse. The resulting scheme is classical per-head Lloyd--Max; the contribution here is the diagnosis of the RoPE incompatibility and the simplest fix that preserves it, not a new quantiser. The fix is method-agnostic: it applies to any RoPE-using backbone and to any KV-eviction scheme that hands off a token-level cache to quantisation; \cref{sec:quant} gives the formal description.


%-----------------------------------------------------------------------%
\subsection{Latency}\label{sec:latency}
%-----------------------------------------------------------------------%

The cumulative-attention term in Stage~1 requires materialising the
post-softmax attention matrix at four sampled layers, which is
nominally incompatible with memory-efficient kernels such as
flash-attention. To verify that this overhead is not observable in
practice, we measure forward-pass wall-clock time at fixed cache
budget $k{=}300$ on five backbones (Qwen~14B/32B, Llama~8B, Mistral~7B
and Mistral-Small~24B) on 18 multi-window SQuAD instances, decomposing
total time into the selection step (\emph{select}) and the
generation step using the selected cache (\emph{generate}).

\begin{table}[t]
\caption{End-to-end forward-pass latency (seconds per query, mean over 18 samples) at $k{=}300$. \emph{Select} is the time to compute the cache-selection scores and gather the retained tokens; \emph{Generate} is the time to produce the answer using the selected cache. \textcolor{red}{\textbf{[PLACEHOLDER]}} StreamingLLM, H$_2$O and SnapKV latency columns relied on the same hand-rolled implementations that this revision withdraws from the accuracy tables; the AstroNet S1 column stands. Faithful re-measurements with upstream baseline implementations are pending.}
\label{tab:latency}
\centering
\small
\setlength{\tabcolsep}{4pt}
\begin{tabular}{l rrr rrr rrr rrr}
\toprule
& \multicolumn{3}{c}{\textbf{StreamingLLM}} & \multicolumn{3}{c}{\textbf{H${}_2$O}} & \multicolumn{3}{c}{\textbf{SnapKV}} & \multicolumn{3}{c}{\textbf{AstroNet S1}} \\
\cmidrule(lr){2-4}\cmidrule(lr){5-7}\cmidrule(lr){8-10}\cmidrule(lr){11-13}
\textbf{Model} & Sel & Gen & Tot & Sel & Gen & Tot & Sel & Gen & Tot & Sel & Gen & Tot \\
\midrule
Llama 3.1 8B      & \textcolor{red}{[P]}&\textcolor{red}{[P]}&\textcolor{red}{[P]} & \textcolor{red}{[P]}&\textcolor{red}{[P]}&\textcolor{red}{[P]} & \textcolor{red}{[P]}&\textcolor{red}{[P]}&\textcolor{red}{[P]} & 0.08 & 0.76 & 1.26 \\
Mistral 7B        & \textcolor{red}{[P]}&\textcolor{red}{[P]}&\textcolor{red}{[P]} & \textcolor{red}{[P]}&\textcolor{red}{[P]}&\textcolor{red}{[P]} & \textcolor{red}{[P]}&\textcolor{red}{[P]}&\textcolor{red}{[P]} & 0.08 & 0.92 & 1.45 \\
Qwen 2.5 14B      & \textcolor{red}{[P]}&\textcolor{red}{[P]}&\textcolor{red}{[P]} & \textcolor{red}{[P]}&\textcolor{red}{[P]}&\textcolor{red}{[P]} & \textcolor{red}{[P]}&\textcolor{red}{[P]}&\textcolor{red}{[P]} & 0.16 & 0.51 & 1.50 \\
Qwen 2.5 32B      & \textcolor{red}{[P]}&\textcolor{red}{[P]}&\textcolor{red}{[P]} & \textcolor{red}{[P]}&\textcolor{red}{[P]}&\textcolor{red}{[P]} & \textcolor{red}{[P]}&\textcolor{red}{[P]}&\textcolor{red}{[P]} & 0.36 & 0.85 & 3.00 \\
Mistral-Small 24B & \textcolor{red}{[P]}&\textcolor{red}{[P]}&\textcolor{red}{[P]} & \textcolor{red}{[P]}&\textcolor{red}{[P]}&\textcolor{red}{[P]} & \textcolor{red}{[P]}&\textcolor{red}{[P]}&\textcolor{red}{[P]} & 0.26 & 1.26 & 2.82 \\
\bottomrule
\end{tabular}
\end{table}

\textcolor{red}{\textbf{[PROSE PLACEHOLDER]}} The Stage~1 selection cost
in isolation is $0.08$--$0.36$~s depending on backbone size, dominated
by a single forward pass on the query window. The generate step
dominates the total per-query cost on every backbone. Comparative
claims against StreamingLLM/H$_2$O/SnapKV are withdrawn pending
faithful re-measurements with the upstream baseline implementations.


%-----------------------------------------------------------------------%
\subsection{Stage~2 contribution transfers across selectors (preliminary)}\label{sec:e3}
%-----------------------------------------------------------------------%

The preceding tables are written around the multiplicative Stage~1 selector specifically. If the headline claim of this paper is that the small learned summary module is an additive plug-in, the contribution must persist when Stage~1's scoring rule is replaced by a different heavy-hitter selector. We report a preliminary three-backbone test of this transferability here; the full six-backbone version will appear once the upstream faithful-baseline queue completes.

The test holds the trained Stage~2 checkpoint fixed and varies only the rule that picks the 284 real-token indices. Four conditions are evaluated on the same SQuAD position-robust samples ($n{=}100$, 4 fact positions, seed 42, $k{=}300$):
\textbf{pure\_S1} (multiplicative selector alone, no Stage~2);
\textbf{pure\_swap} (SnapKV-style scoring: max-pool kernel-7 instead of avg-pool kernel-5, recent strip from the last 32 tokens of cumulative; no Stage~2);
\textbf{hybrid\_S1} (multiplicative selector + 16 Stage~2 tokens, the published AstroNet configuration);
\textbf{hybrid\_swap} (SnapKV-style scoring + 16 Stage~2 tokens). The contribution of Stage~2 is then read off as $\Delta_{\text{S1}} = $ hybrid\_S1 $-$ pure\_S1 and $\Delta_{\text{swap}} = $ hybrid\_swap $-$ pure\_swap.

\begin{table}[t]
\caption{Stage~2 transferability across selectors (preliminary, three backbones). Same trained Stage~2 checkpoint, two alternative selectors for the 284 real-token slots, $k{=}300$ total. SQuAD position-robust mean across 4 fact positions, $n{=}100$ samples per position, seed 42. $\Delta_{\text{S1}}$ is the Stage~2 contribution paired with the multiplicative selector; $\Delta_{\text{swap}}$ is the contribution paired with SnapKV-style scoring under the same protocol. Equal $\Delta$ across selectors indicates Stage~2 is an additive plug-in rather than a co-adapted partner to one specific selector. Remaining three backbones (Qwen 32B, Mistral 7B, Mistral-Small 24B) pending.}
\label{tab:e3_transfer}
\centering
\small
\setlength{\tabcolsep}{5pt}
\begin{tabular}{l rrr@{\hskip 8pt} rrr}
\toprule
& \multicolumn{3}{c}{\textbf{Multiplicative selector}} & \multicolumn{3}{c}{\textbf{SnapKV-style selector}} \\
\cmidrule(lr){2-4} \cmidrule(lr){5-7}
\textbf{Model} & pure\_S1 & hybrid\_S1 & $\Delta_{\text{S1}}$ & pure\_swap & hybrid\_swap & $\Delta_{\text{swap}}$ \\
\midrule
Qwen 2.5 7B   & 57.8\% & 62.5\% & $+4.7$ & 62.0\% & 65.2\% & $+3.2$ \\
Llama 3.1 8B  & 50.0\% & 57.5\% & $+7.5$ & 51.0\% & 58.2\% & $+7.3$ \\
Qwen 2.5 14B  & 62.5\% & 65.8\% & $+3.2$ & 66.2\% & 73.0\% & $+6.8$ \\
\midrule
\textbf{Mean (3 models)} & 56.8\% & 61.9\% & $\boldsymbol{+5.2}$ & 59.7\% & 65.5\% & $\boldsymbol{+5.7}$ \\
\bottomrule
\end{tabular}
\end{table}

On all three backbones tested the Stage~2 contribution is positive under both selectors and the two $\Delta$ values are within $\pm 4$ pp of each other ($\Delta_{\text{S1}}$ range $+3.2$ to $+7.5$, $\Delta_{\text{swap}}$ range $+3.2$ to $+7.3$). Averaged across the three backbones, $\Delta_{\text{S1}} = +5.2$ pp and $\Delta_{\text{swap}} = +5.7$ pp; the two means are statistically indistinguishable at this sample size. The reading is that Stage~2 is not co-adapted to Stage~1's specific scoring rule: it transfers as an additive plug-in to a different heavy-hitter selector under the same operating protocol. \textit{The full six-backbone version of this test, plus the matched-budget control ($k{=}284$), is in progress and will be substituted here when the queue completes.}


%-----------------------------------------------------------------------%
\subsection{Ablations}\label{sec:ablations}
%-----------------------------------------------------------------------%

We report three ablations that bear directly on the headline claims: the Stage~2 contribution under a matched real-token budget, the functional form of the cross-window aggregator, and the size of the learned summary.

The first ablation isolates the Stage~2 contribution from any potential budget effect. The selection-plus-summary pipeline differs from pure selection in two ways: it allocates 16 of its 300 slots to learned summary tokens, and it injects content that is not present in the original cache. To separate these effects we compare three conditions at $n{=}20$ windows: pure Stage~1 at the full budget ($k{=}300$ real tokens), pure Stage~1 at the matched budget ($k{=}284$ real tokens, the same real-token count the hybrid uses), and Stage~1+2 ($16{+}284$ slots). \Cref{tab:kconfound} shows that the matched-budget Stage~1 sometimes outperforms unmatched Stage~1 (the $k{=}284$ baseline removes noisy borderline tokens), so the matched control is the conservative diagnostic.

\begin{figure}[t]
\centering
\includegraphics[width=0.85\textwidth]{figs/fig_kconfound_dumbbell.pdf}
\caption{\textbf{Matched-budget Stage~2 contribution} (companion to Table~\ref{tab:kconfound}). Each row plots S1 at $k{=}284$ (matched real-token budget) and S1+S2 ($16{+}284$) as a dumbbell; the right annotation is the multi-seed $\Delta$ ($\pm$ std). All six backbones show a positive $\Delta$; the Mistral-Small~24B value is from the longer-context retrain (the original SQuAD-only training produced $\Delta=-16$).}
\label{fig:kconfound}
\end{figure}

\begin{table}[t]
\caption{Stage~2 contribution at matched budget. Accuracy (\%) on needle retrieval at $n{=}20$ windows. The $\Delta$ column reports Stage~1+2 minus Stage~1 at the matched real-token budget of $k{=}284$, isolating the contribution of the 16 learned tokens specifically. The Mistral-Small~24B row uses the longer-context retrained checkpoint discussed in Section~\ref{sec:discussion}; the original checkpoint produced $\Delta=-16$\,pp on this diagnostic before retraining.}
\label{tab:kconfound}
\centering
\small
\setlength{\tabcolsep}{6pt}
\begin{tabular}{l c c c c}
\toprule
\textbf{Model} & \textbf{S1, $k{=}300$} & \textbf{S1, $k{=}284$} & \textbf{S1+S2 (16+284)} & $\boldsymbol{\Delta}$ \\
\midrule
Qwen 2.5 7B       & 74 & 82 & 88  & $+6$  \\
Qwen 2.5 14B      & 78 & 80 & 100 & $+20$ \\
Qwen 2.5 32B      & 48 & 60 & 100 & $+40$ \\
Llama 3.1 8B      & 60 & 40 & 92  & $+52$ \\
Mistral 7B        &  8 & 10 & 32  & $+22$ \\
Mistral-Small 24B & 72 & 72 & 81  & $+9$  \\
\bottomrule
\end{tabular}
\end{table}

The matched-budget contribution of the 16 learned tokens is positive on all six backbones, ranging from $+6$\,pp on Qwen~7B to $+52$\,pp on Llama~8B; the Mistral-Small~24B value of $+9$\,pp is from the longer-context retrain that resolved a $-16$\,pp regression observed with the original $n{=}5$ SQuAD-only training. The diagnostic confirms that the hybrid gain is not a budget artefact: a Stage~1 pipeline at the same 284-real-token budget would not reach the Stage~1+2 numbers.

The second ablation probes the functional form of the cross-window aggregator. The Stage~2 summary maintains a state $\mathbf{g}$ through an exponential moving average with learnable rate $\alpha$ (\cref{sec:s2}). We replace it at evaluation time with three alternatives: $\alpha{=}1$ (last window only, no accumulation), $\alpha{=}0$ (frozen initial state, no sensing), and uniform mean over all sensed windows. \Cref{tab:bio} reports the resulting accuracy on the $n{=}5$-window needle task (50 trials) for each backbone.

\begin{table}[t]
\caption{Effect of the Stage~2 state-update rule on $n{=}5$-window needle accuracy (\%, 50 trials). \emph{Default}: learnable EMA. $\alpha{=}1$: last-window-only. $\alpha{=}0$: frozen initial state. \emph{Mean}: uniform average over windows.}
\label{tab:bio}
\centering
\small
\setlength{\tabcolsep}{6pt}
\begin{tabular}{l c c c c}
\toprule
\textbf{Model} & \textbf{Default (EMA)} & $\boldsymbol{\alpha{=}1}$ & $\boldsymbol{\alpha{=}0}$ & \textbf{Mean} \\
\midrule
Qwen 2.5 7B       & 90  & 90  & 38 & 90  \\
Qwen 2.5 14B      & 100 & 98  & 84 & 100 \\
Qwen 2.5 32B      & 100 & 100 & 40 & 100 \\
Llama 3.1 8B      & 90  & 90  & 44 & 90  \\
Mistral 7B        & 32  & 32  & 10 & 32  \\
Mistral-Small 24B & 80  & 82  & 64 & 80  \\
\bottomrule
\end{tabular}
\end{table}

On all six backbones the EMA default, the last-window-only setting, and the uniform mean give equivalent accuracy; the only condition that consistently degrades is $\alpha{=}0$ (the summary is starved of any window-specific content and reduces to a learned bias). The interpretation is methodological rather than architectural: the specific cross-window aggregation rule of Stage~2 is not load-bearing -- any non-trivial aggregator that reads the per-window sensed signal works equally well, including the simplest possible choice of a uniform mean. The persistent state itself, however, is necessary: removing it ($\alpha{=}0$) loses 16--52 percentage points across the six backbones. We retain the learnable EMA as the canonical form because it is the natural training-time parameterisation, but the ablation supports a simplification toward a uniform mean if a future deployment prefers a parameter-free aggregator.

The third ablation sweeps the number of learned summary tokens $K$ within a fixed total budget of 300, replacing $K$ real tokens with $K$ summary tokens. On Qwen~7B SQuAD, accuracy is essentially flat across $K{\in}\{4, 8, 16\}$ (83, 84, 83\%), with a single percentage-point degradation at $K{=}32$ (82\%). The architecture is therefore robust to the precise count of summary tokens: four are already sufficient on this controlled benchmark, sixteen is on the same plateau, and beyond thirty the displacement of real tokens begins to dominate. We retain $K{=}16$ for the canonical configuration because it lies on the plateau, sits comfortably below the displacement regime, and matches the configuration trained for all six backbones.

A fourth ablation probes the sensitivity of the multiplicative cross-window score to the question text itself. \Cref{fig:queryablation} shows hybrid accuracy across three substitutions of the query-window content during the selection step: \emph{Real} (the canonical question text), \emph{Empty} (a literal empty string), and \emph{Trailing-32} (the last 32 tokens of the last fact window in place of the question). The Qwen and Llama backbones show parallel downward slopes from Real to Empty to Trailing-32 of 13--24~pp; the two Mistral backbones flatten visibly (3~pp on Mistral~7B and 6~pp on Mistral-Small~24B), consistent with weaker question-conditioned selection on those backbones. The graceful Real$\to$Empty$\to$Trailing-32 degradation supports the interpretation that the multiplicative score is genuinely question-conditioned rather than reading off a lexical artefact of the input format.

\begin{figure}[t]
\centering
\includegraphics[width=0.85\textwidth]{figs/fig_queryablation_slope.pdf}
\caption{\textbf{Question-text sensitivity of the multiplicative selection score.} Hybrid accuracy ($n{=}100$ SQuAD samples, mean over 4 fact positions, seed 42) when the query window text is replaced during selection with three alternatives: the real question, the empty string, and the last 32 tokens of the last fact window. Parallel downward slopes on the Qwen and Llama backbones indicate genuine question-conditioned selection; the Mistral lines flatten, consistent with weaker question-conditioning on those backbones.}
\label{fig:queryablation}
\end{figure}
